\begin{Remark}
\label{rem:chap1-angular-momentum-redundant}
既然已经证明应力张量对称, 线动量演化方程显然蕴含角动量演化方程. 因此从现在起, 不必再把后一个方程加入控制流体运动的方程组.

利用应力张量的对称性, 能量演化方程\eqref{eq:chap1-energy-local}可写成
\begin{equation}
\label{eq:chap1-energy-conservative-symmetric}
  \frac{\partial(\rho E)}{\partial t}+\operatorname{div}(\rho E v)
  -\operatorname{div}(\sigma\cdot v)+\operatorname{div}\phi=\rho f\cdot v.
\end{equation}
\end{Remark}

\begin{Remark}
\label{rem:chap1-conservative-nonconservative}
目前得到的式\eqref{eq:chap1-linear-momentum-local}和\eqref{eq:chap1-energy-conservative-symmetric}称为守恒形式, 因为它们写成 $(t,x)$ 中某个量的散度等于源项的形式. 对形式计算或物理解读而言, 有时更适合使用下面的非守恒形式.

线动量方程:
\begin{equation}
\label{eq:chap1-linear-momentum-nonconservative}
  \rho\left(\frac{\partial v}{\partial t}+(v\cdot\nabla)v\right)-\operatorname{div}\sigma=\rho f.
\end{equation}
能量方程:
\begin{equation}
\label{eq:chap1-energy-nonconservative}
  \rho\left(\frac{\partial E}{\partial t}+(v\cdot\nabla)E\right)
  -\operatorname{div}(\sigma\cdot v)+\operatorname{div}\phi=\rho f\cdot v.
\end{equation}
这些方程从前面的方程出发, 利用质量平衡方程\eqref{eq:chap1-mass-balance}以及\MBUnresolvedReference{eq:app-div-product}{Formulae~(A.4), (A.10)}作代数运算得到. 从形式观点看, 两种写法完全等价. 从数学观点看, 当处理正则性不高的弱解时, 两种表述是否等价, 甚至非守恒形式是否有意义, 并不总是清楚的.
\end{Remark}

\subsection{基本定律: Newtonian 流体与热力学定律}
\label{sec:chap1-fundamental-laws}

上一节给出了质量守恒, 线动量和能量演化方程的一般形式, 用以描述受体力作用且所研究区域内部没有物质源或热源的流体流动. 这些方程仍包含若干尚未确定的量: 应力张量 $\sigma$(目前只知道它对称), 热通量 $\phi$ 和比内能 $e$. 后两个量通过热力学定律与描述流动的其他场相联系, 我们将在 Section~\ref{subsec:chap1-internal-energy} 回到这些量.

\subsubsection{静止流体}
\label{subsec:chap1-fluids-at-rest}

\begin{Definition}
\label{def:chap1-hydrostatic-pressure}
在静止流体中, 作用于流体元表面微元上的应力方向与该表面的外法向量方向相反. 此外, 应力的模与方向无关. 其模记为 $p$, 称为流体的静水压力.
\end{Definition}

由此定义, 对静止流体的每一点都有 $\tau(\nu)=-p\nu$, 其中 $p\geq0$. 因而应力张量写为 $\sigma=-p\operatorname{Id}$. 当流体运动时, 将应力张量表示成
\[
  \sigma=T-p\operatorname{Id},
\]
以分离压力效应和运动效应. 新张量 $T$ 显然对称, 称为黏性应力张量. 这里假定静水压力 $p$ 与\MBUnresolvedReference{app:thermodynamics-supplement}{Appendix~B} 中由材料分子状态定义的热力学压力 $P$ 完全一致.

\subsubsection{Newton 假设}
\label{subsec:chap1-newton-hypothesis}

考虑黏性流体的平面 Couette 流, 另见 Section~\ref{subsec:chap1-couette-flow}. 流体位于两块假定为无限大的平行平板之间. 在大小恒定且沿 $e_1$ 方向的应力 $Fe_1$ 作用下, 上板以速度 $U_0e_1$ 运动, 下板保持静止. 对水等经典流体, 实验表明使平板以速度 $U_0$ 运动所需的黏性应力为
\begin{equation}
\label{eq:chap1-couette-stress}
  F=\mu\frac{U_0}{d},
\end{equation}
其中常数 $\mu$ 仅依赖于所考察的流体. 换言之, 应力 $F$ 与流动中的速度梯度 $U_0/d$ 成正比.

这种应力效应是流体对剪切的一种阻力. 若停止施加驱动力 $F$, 流体经过一段时间后将恢复静止. 例如水的速度会随时间衰减. 这样的流体称为黏性流体. 黏性效应在流体发生形变时出现. 对 Couette 流可直观地说, 高度 $x_3=C+h$ 处粒子的水平速度大于高度 $x_3=C$ 处粒子的速度, 从而产生流体形变.

\begin{Definition}
\label{def:chap1-rigid-body-behaviour}
考虑速度场为 $v$ 的流动. 任取三个彼此很近的点 $x_0$, $x_1=x_0+(\delta x)\xi_1$ 和 $x_2=x_0+(\delta x)\xi_2$. 记 $\widetilde x_0$, $\widetilde x_1$, $\widetilde x_2$ 为时刻 $t$ 分别位于 $x_0$, $x_1$, $x_2$ 的流体粒子在时刻 $t+\delta t$ 的位置.

若对方向 $\xi_1,\xi_2$ 的任意选择, 三角形 $x_0x_1x_2$ 与 $\widetilde x_0\widetilde x_1\widetilde x_2$ 在关于 $\delta t$ 和 $\delta x$ 的一阶精度内可以重合, 则称流体在时刻 $t$ 点 $x_0$ 表现为刚体. 更确切地说, 对所有 $i,j,k$,
\begin{equation}
\label{eq:chap1-rigid-body-inner-product}
  \overrightarrow{\widetilde x_i\widetilde x_j}\cdot\overrightarrow{\widetilde x_i\widetilde x_k}
  =\bigl(\overrightarrow{x_ix_j}\cdot\overrightarrow{x_ix_k}\bigr)
  \left(1+O(\delta t^2+\delta x^2)\right).
\end{equation}
\end{Definition}

\begin{Proposition}[Definition and Proposition]
\label{prop:chap1-strain-rate}
流体在 $(t,x)$ 表现为刚体, 当且仅当张量
\[
  D(v)=\frac12(\nabla v+{}^t\!\nabla v)
  =\left(\frac12\left(\frac{\partial v_i}{\partial x_j}+\frac{\partial v_j}{\partial x_i}\right)\right)_{1\leq i,j\leq3}
\]
在 $(t,x)$ 为零. 张量 $D(v)$ 称为流动的应变率张量.
\end{Proposition}

\begin{Proof}
记 $s\mapsto x_i(s)$ 为时刻 $t$ 从 $x_i$ 出发的轨迹, 因而 $\widetilde x_i=x_i(t+\delta t)$. 由特征方程\eqref{eq:chap1-characteristic-equation},
\[
  x_i'(s)=v(s,x_i(s)),\qquad
  x_i''(s)=\left(\frac{\partial v}{\partial t}+(v\cdot\nabla)v\right)(s,x_i(s)).
\]
对每个 $j\in\{0,1,2\}$ 作 Taylor 展开, 并利用特征方程的解对初值的连续依赖以及速度场的正则性, 得
\[
  \widetilde x_j-\widetilde x_i=x_j-x_i+\delta t\bigl(v(t,x_j)-v(t,x_i)\bigr)+O(\delta t^2\delta x).
\]
再对 $v(t,x_j)-v(t,x_i)$ 作一阶展开, 得
\[
\overrightarrow{\widetilde x_0\widetilde x_1}
=\overrightarrow{x_0x_1}+\delta t\delta x\nabla v(t,x_0)\cdot\xi_1+O(\delta t\delta x^2)+O(\delta t^2\delta x),
\]
以及对 $\xi_2$ 的相同公式. 取内积后得到
\begin{align*}
\overrightarrow{\widetilde x_0\widetilde x_1}\cdot\overrightarrow{\widetilde x_0\widetilde x_2}
&=\overrightarrow{x_0x_1}\cdot\overrightarrow{x_0x_2}\\
&\quad+\delta t\delta x^2\bigl((\nabla v(t,x_0)\xi_1)\cdot\xi_2+(\nabla v(t,x_0)\xi_2)\cdot\xi_1\bigr)\\
&\quad+O(\delta t\delta x^3)+O(\delta t^2\delta x^2).
\end{align*}
因此条件\eqref{eq:chap1-rigid-body-inner-product}等价于对所有 $\xi_1,\xi_2$ 有
\[
  (\nabla v(t,x_0)\xi_1)\cdot\xi_2+(\nabla v(t,x_0)\xi_2)\cdot\xi_1=0.
\]
这恰好说明 $\nabla v(t,x_0)$ 反对称, 等价地, $D(v)$ 在 $(t,x_0)$ 为零.
\end{Proof}

\begin{Remark}
\label{rem:chap1-rigid-velocity-field}
可以证明, 见\MBUnresolvedReference{lem:chap4-rigid-velocity}{Lemma~IV.7.5}, 若速度场 $v$ 对所有 $(t,x)$ 满足 $D(v)=0$, 则必有
\[
  v(t,x)=v_0(t)+\omega(t)\wedge r,
\]
其中 $r$ 是固定参考系中的位置向量. 这意味着运动是空间上一致速度的平移与绕固定轴旋转的叠加, 正是刚体可能具有的运动.
\end{Remark}

\begin{Remark}
\label{rem:chap1-strain-rate-identities}
直接计算给出
\[
\operatorname{Tr}(D(v))=\operatorname{Tr}(\nabla v)=\operatorname{div}v,
\quad \operatorname{div}({}^t\!\nabla v)=\nabla(\operatorname{div}v),
\]
\[
\operatorname{div}(2D(v))=\Delta v+\nabla(\operatorname{div}v),
\quad \operatorname{div}((\operatorname{div}v)\operatorname{Id})=\nabla(\operatorname{div}v).
\]
这些公式将在后文使用.
\end{Remark}

流体形变率的全部信息都包含在速度梯度张量 $\nabla v$ 的对称部分 $D(v)$ 中. 流体的标准黏性行为概括为: 黏性应力张量 $T$ 仅依赖于 $D(v)$; $T$ 对 $D(v)$ 的依赖是线性的; 连接 $T$ 与 $D(v)$ 的关系是各向同性的. 实验上满足这些性质的流体称为 Newtonian 流体. 在温度, 压力或速度的极端条件下, Newtonian 流体也可能表现出非 Newtonian 行为.

第一条假设来自前面的考虑. 线性假设是式\eqref{eq:chap1-couette-stress}的推广, 可写为 $T=L(D(v))$, 其中 $L$ 是对称张量空间到自身的线性映射. 各向同性表示改变标准正交参考系时流体的流变性质不变. 其确切含义是: 对所有正交矩阵 $P$ 和所有对称张量 $d$,
\begin{equation}
\label{eq:chap1-isotropic-viscosity}
  L({}^t\!P dP)={}^t\!P L(d)P.
\end{equation}

\begin{Proposition}
\label{prop:chap1-newtonian-stress-law}
对 Newtonian 流体, 黏性应力张量关于应变率张量的定律必具有形式
\[
  T=2\mu D(v)+\lambda(\operatorname{div}v)\operatorname{Id},
\]
其中 $\lambda$ 和 $\mu$ 是实系数.
\end{Proposition}

\begin{Proof}
令 $d$ 为对称张量, $x$ 为 $d$ 对应特征值 $\Lambda$ 的特征向量. 取关于垂直于 $x$ 的超平面的对称变换矩阵 $P$, 即 $Px=-x$, 且 $P$ 在 $x$ 的正交补上的限制为恒等映射. 由 ${}^t\!PdP=d$ 和式\eqref{eq:chap1-isotropic-viscosity}, 得 $L(d)P=PL(d)$. 将其作用于 $x$, 可知 $L(d)x$ 与 $x$ 共线, 所以 $x$ 也是 $L(d)$ 的特征向量.

记 $E_i$ 为除 $(i,i)$ 元为 $1$ 外其余系数均为零的矩阵. 存在 $\lambda,\mu$ 使 $L(E_1)=2\mu E_1+\lambda\operatorname{Id}$. 令 $Q_i$ 为交换指标 $1$ 与 $i$ 的置换矩阵, 则 $E_i={}^t\!Q_iE_1Q_i$. 再用式\eqref{eq:chap1-isotropic-viscosity}, 得 $L(E_i)=2\mu E_i+\lambda\operatorname{Id}$. 因而对任意对角张量 $d$,
\[
  L(d)=2\mu d+\lambda(\operatorname{Tr}d)\operatorname{Id}.
\]
对任意对称张量 $d$, 取正交矩阵 $Q$ 使 ${}^t\!QdQ$ 对角, 再次应用各向同性即得同一公式. 最后利用 $\operatorname{Tr}(D(v))=\operatorname{div}v$, 得所述结论.
\end{Proof}

Newtonian 流体的应力张量一般形式为
\[
  \sigma=T-p\operatorname{Id}=2\mu D(v)+(\lambda\operatorname{div}v-p)\operatorname{Id}.
\]

\subsubsection{热力学第二定律的推论}
\label{subsec:chap1-second-law}

热力学第二定律意味着黏性应力必然是耗散的, 即这些力的功率必须非负. 利用\MBUnresolvedReference{app:thermodynamics-supplement}{Appendix~B} 中回顾的结果, 局部比熵方程可写成守恒形式
\[
  \frac{\partial(\rho s)}{\partial t}+\operatorname{div}(\rho sv)
  +\operatorname{div}\left(\frac{\phi}{T}\right)-\frac{\operatorname{Tr}(T D(v))}{T}=0,
\]
其中 $s$ 是流动中的比熵, $T$ 是温度. 对任意随流动演化的流体元 $(\Omega_t)_t$, 输运定理给出总熵演化. 在时刻 $t$ 与 $t+\delta t$ 之间积分, 并令 $S(t)=\int_{\Omega_t}\rho s\,\mathrm{d}X$, 得
\begin{equation}
\label{eq:chap1-clausius-balance}
\begin{aligned}
S(t+\delta t)-S(t)-\int_t^{t+\delta t}\int_{\partial\Omega_s}\frac{-\phi\cdot\nu}{T}\,\mathrm{d}y\,\mathrm{d}s
&=\int_t^{t+\delta t}\int_{\Omega_s}\frac{-\phi\cdot\nabla T}{T^2}\,\mathrm{d}X\,\mathrm{d}s\\
&\quad+\int_t^{t+\delta t}\int_{\Omega_s}\frac{\operatorname{Tr}(T D(v))}{T}\,\mathrm{d}X\,\mathrm{d}s.
\end{aligned}
\end{equation}
左端为 $\Delta S-\Delta Q/T$, 其中 $\Delta Q$ 是该流体元在两个时刻之间通过边界获得的热量. Clausius 定理指出此量非负, 若过程不可逆则为正. 由此与式\eqref{eq:chap1-clausius-balance}推出
\begin{equation}
\label{eq:chap1-heat-flux-dissipation}
  \phi\cdot\nabla T\leq0,\qquad\text{对任意可能的 $\nabla T$},
\end{equation}
以及
\begin{equation}
\label{eq:chap1-viscous-dissipation}
  \operatorname{Tr}(T D(v))\geq0,\qquad\text{对任意可能的 $D(v)$}.
\end{equation}

用上一段记号, 式\eqref{eq:chap1-viscous-dissipation}表示对称张量空间上的二次型 $d\mapsto P(d)=L(d):d$ 必须半正定. 由各向同性, 只需研究对角张量. 在向量 $(1,-1,0)^t$, $(1,0,-1)^t$, $(1,1,1)^t$ 构成的基中, 相应矩阵为
\[
  \begin{pmatrix}4\mu&2\mu&0\\2\mu&4\mu&0\\0&0&6\mu+9\lambda\end{pmatrix}.
\]
因此半正定当且仅当
\begin{equation}
\label{eq:chap1-viscosity-conditions}
  \mu\geq0,\qquad 2\mu+3\lambda\geq0.
\end{equation}
系数 $\mu$ 称为动力黏度, $2\mu/3+\lambda$ 称为体积黏度. 在许多情形下体积黏度很小, 可忽略, 这就是 Stokes 假设:
\begin{equation}
\label{eq:chap1-stokes-hypothesis}
  \lambda=-\frac23\mu.
\end{equation}
此时
\[
  T=2\mu\left(D(v)-\frac13(\operatorname{div}v)\operatorname{Id}\right),
\]
且其迹为零, 所有应力的各向同性分量都包含在压力 $p$ 中. 若动力黏度 $\mu$ 很小且可忽略, 令 $\mu=0$, 流体称为理想流体或无黏流体.

最后给出热通量 $\phi$ 的一般表达式. 若假定 $\phi$ 只依赖于温度及其梯度, 且对 $\nabla T$ 的依赖在惯性参考系变换下不变, 则 Fourier 定律指出 $\phi$ 与温度梯度成正比. 式\eqref{eq:chap1-heat-flux-dissipation}进一步说明二者方向相反, 因而
\begin{equation}
\label{eq:chap1-fourier-law}
  \phi=-k\nabla T,
\end{equation}
其中 $k\geq0$ 称为流体的热导率, 它可以依赖于温度或压力等流体特性.

\subsubsection{比内能方程}
\label{subsec:chap1-internal-energy}

总能量与比内能满足 $E=e+\tfrac12|v|^2$. 将非守恒动量方程\eqref{eq:chap1-linear-momentum-nonconservative}与 $v$ 作内积, 得到动能演化方程. 从能量方程\eqref{eq:chap1-energy-nonconservative}中减去它, 并利用 $\sigma$ 的对称性以及 $\sigma=T-p\operatorname{Id}$, 得
\[
  \rho\left(\frac{\partial e}{\partial t}+v\cdot\nabla e\right)
  +\operatorname{div}\phi+p\operatorname{div}v-\operatorname{Tr}(T\nabla v)=0.
\]
对 Newtonian 流体,
\[
\operatorname{Tr}(T\nabla v)=\operatorname{Tr}(TD(v))
=\lambda(\operatorname{div}v)^2+2\mu D(v):D(v)
=\lambda(\operatorname{div}v)^2+2\mu|D(v)|^2.
\]
在条件\eqref{eq:chap1-viscosity-conditions}下, 此量总是非负, 与热力学第二定律一致. 应用 Fourier 定律\eqref{eq:chap1-fourier-law}, 比内能的演化方程为
\begin{equation}
\label{eq:chap1-internal-energy}
\rho\left(\frac{\partial e}{\partial t}+v\cdot\nabla e\right)
-\operatorname{div}(k\nabla T)+p\operatorname{div}v
-\lambda(\operatorname{div}v)^2-2\mu|D(v)|^2=0.
\end{equation}
其守恒形式为
\[
\frac{\partial(\rho e)}{\partial t}+\operatorname{div}(\rho ev)
-\operatorname{div}(k\nabla T)+p\operatorname{div}v
-\lambda(\operatorname{div}v)^2-2\mu|D(v)|^2=0.
\]

\subsubsection{熵与温度表述}
\label{subsec:chap1-entropy-temperature}

利用\MBUnresolvedReference{app:thermodynamics-supplement}{Appendix~B} 中的热力学概念, 可以用比熵和温度作为变量写出等价的守恒方程. 假定流体粒子内部结构的弛豫时间远短于流体整体运动的特征时间, 因而流体在每一点和每一时刻均处于热力学平衡.

将 $(s,\rho)\mapsto e(s,\rho)$ 的微分写为
\[
  \mathrm{d}e=T\,\mathrm{d}s-p\,\mathrm{d}\left(\frac1\rho\right).
\]
在 Lagrangian 层次使用物质导数并利用质量守恒方程\eqref{eq:chap1-mass-balance}, 得
\[
  \frac{De}{Dt}=T\frac{Ds}{Dt}+\frac{p}{\rho^2}\frac{D\rho}{Dt}
  =T\frac{Ds}{Dt}-\frac p\rho\operatorname{div}v.
\]
所以式\eqref{eq:chap1-internal-energy}成为比熵方程
\[
  \rho T\frac{Ds}{Dt}-\operatorname{div}(k\nabla T)
  -\lambda(\operatorname{div}v)^2-2\mu|D(v)|^2=0.
\]

为得到温度方程, 使用比内能关于温度和比容的微分式
\[
  \mathrm{d}e=c_v\,\mathrm{d}T+\left(T\left(\frac{\partial p}{\partial T}\right)_\rho-p\right)\mathrm{d}\left(\frac1\rho\right).
\]
再次利用质量守恒方程,
\[
  \frac{De}{Dt}=c_v\frac{DT}{Dt}+\frac1\rho T\left(\frac{\partial p}{\partial T}\right)_\rho\operatorname{div}v-p\frac1\rho\operatorname{div}v.
\]
因此温度形式为
\[
\rho c_v\frac{DT}{Dt}+T\left(\frac{\partial p}{\partial T}\right)_\rho\operatorname{div}v
-\operatorname{div}(k\nabla T)-\lambda(\operatorname{div}v)^2-2\mu|D(v)|^2=0.
\]
若 $c_v$ 为常数, 其守恒形式为
\[
\frac{\partial(\rho c_vT)}{\partial t}+\operatorname{div}(\rho c_vTv)
+T\left(\frac{\partial p}{\partial T}\right)_\rho\operatorname{div}v
-\operatorname{div}(k\nabla T)-\lambda(\operatorname{div}v)^2-2\mu|D(v)|^2=0.
\]
